% ============================================================================
% CUT CONTENT — preserved for advisor review
% This file contains material removed from the main paper during the
% research-ready cleanup. Nothing here is included in the compiled paper.
% ============================================================================

% ----------------------------------------------------------------------------
% CUT 1: Linear Projection Transfer (was Section 3.6 and Appendix B)
% Reason: No experimental evidence yet. Data exists on HPC for Experiment B
% (ACT→OpenVLA and reverse projection). Kept as future work note.
% ----------------------------------------------------------------------------

% \subsection{Cross-Model Safety Transfer}
% \label{sec:transfer_cut}
% To test whether the safety manifold is \emph{shared} across architectures, we train a safety head on one policy (e.g., ACT, $m=256$) and evaluate it on a different architecture (e.g., OpenVLA, $m=4096$) through a learned linear projection:
% \begin{equation}
% \mathbf{h}^{\mathrm{proj}}_t = \mathrm{LayerNorm}\bigl(W \mathbf{h}_t + \mathbf{b}\bigr), \quad W \in \mathbb{R}^{m_{\mathrm{target}} \times m_{\mathrm{source}}}.
% \end{equation}
% We measure \textbf{transfer efficiency}: the fraction of the native (retrained) AUC recovered by the projected safety head. Values near 100\% indicate that the safety-relevant directions are linearly aligned across architectures.
%
% Training uses MSE reconstruction loss on the source model's hidden states, optimized with Adam ($\eta=10^{-3}$) for 40 epochs. Transfer efficiency is defined as:
% \begin{equation}
% \eta_{\mathrm{transfer}} = \frac{\mathrm{AUC}_{\mathrm{projected}}}{\mathrm{AUC}_{\mathrm{fresh}}}.
% \end{equation}
%
% NOTE: Multi-model rollout data exists in data/multi_model/ for ACT and OpenVLA.
% Experiment B (~1 GPU-day) would train this projection and measure AUC transfer.

% ----------------------------------------------------------------------------
% CUT 2: Yahboom Bring-Up Checklist (was Appendix and Experimental Setup)
% Reason: No autonomous task-completion results. Infrastructure status is not
% a research contribution. Will be relevant when real robot arm results arrive.
% ----------------------------------------------------------------------------

% \section{Yahboom Embodied Deployment Checklist}
% \begin{table}[h]
% \centering
% \caption{Embodied bring-up status on Yahboom 7-DoF + Jetson Orin.}
% \begin{tabular}{lcc}
% \toprule
% Component & Status & Evidence \\
% \midrule
% SSH + network & \checkmark & Passwordless SSH configured \\
% CUDA + PyTorch & \checkmark & 4-GPU sharding verified \\
% Camera stream & \checkmark & 640$\times$480 @ 30fps \\
% Arm API (status/home/angles) & \checkmark & REST API tested \\
% Gripper open/close & \checkmark & Tested via API \\
% OpenVLA inference (sharded) & \checkmark & $\sim$12s load, $\sim$2s/step \\
% Safety MLP inference & \checkmark & $<$1ms on CUDA \\
% Workspace bounding & \checkmark & $\pm$50mm XY, $\pm$30mm Z \\
% Saturation detection & \checkmark & L$\infty > 0.95$ triggers freeze \\
% Stuck-action termination & \checkmark & 4-step identical $\Rightarrow$ abort \\
% Camera quality diagnostics & \checkmark & Brightness/contrast/focus verified \\
% VLA action-conditioning diagnostics & Partial & Variance acceptable; conditioning under review \\
% Autonomous pick/place benchmark & Pending & Blocked on conditioning fix \\
% Safety steering evaluation & Pending & Requires autonomous baseline \\
% \bottomrule
% \end{tabular}
% \end{table}

% Embodied platform subsection from experimental_setup.tex:
% \subsection{Embodied Platform: Yahboom 7-DoF + Jetson Orin}
% \paragraph{Hardware.} Yahboom 7-DoF arm with parallel gripper, connected via serial
% to a Jetson Orin (JetPack 6.x, 4x GPU, 44GB combined).
% \paragraph{Policy inference.} OpenVLA-7B loaded with device_map=auto across 4 GPUs.
% \paragraph{Safety interlocks.} Workspace bounding, velocity saturation detection,
% repeated-action detection, human-reachable emergency stop.
% \paragraph{Acceptance diagnostics.} vla_control_diagnostics.py verifies action
% standard deviation, prompt-conditioned delta, frame-conditioned delta, gripper variance.

% ----------------------------------------------------------------------------
% CUT 3: Figures 15-20 (qualitative filler)
% Reason: These are descriptive visualizations, not evidence for any claim.
% Keep data artifacts; regenerate if needed for camera-ready with stronger framing.
% ----------------------------------------------------------------------------

% Figure 15: task_safety_phase_space.png — intervention rate vs correction magnitude
% Figure 16: libero_annotated_panels.png — scene snapshots
% Figure 17: latent_pca_success_failure.png — PCA of hidden states
% Figure 18: latent_pca_task_colored.png — PCA colored by task
% Figure 19: rollout_frame_metric_storyboard.png — frame timeline
% Figure 20: steering_task_dashboard.png — per-task dashboard

% ----------------------------------------------------------------------------
% CUT 4: "Non-overlap statement" paragraph from related_work.tex
% Reason: Defensive framing that weakens the paper. The contributions speak
% for themselves without explicitly disclaiming overlap.
% ----------------------------------------------------------------------------

% \paragraph{Non-overlap statement with SAFE-style work.}
% Our method should be read as complementary to SAFE-style detection: we reuse
% failure prediction signals but add a bounded correction controller and evaluate
% intervention behavior (rate, correction magnitude, trajectory deviation) in
% addition to binary success. We do \emph{not} claim to replace anomaly detection;
% the novelty claim is the post-hoc latent correction mechanism and its
% OOD-oriented evaluation protocol.

% ----------------------------------------------------------------------------
% CUT 5: "Transparent statistical reporting" as contribution
% Reason: Good practice is not a contribution.
% ----------------------------------------------------------------------------

% \item[\textbf{C2}] \textbf{Reproducible statistical pipeline.} We provide a
% frozen-data evaluation workflow with pooled and paired tests, Wilson intervals,
% and multiplicity control for pre-specified contrasts.

% ----------------------------------------------------------------------------
% CUT 6: Embodied results subsection from results.tex
% Reason: No autonomous task-completion data. Bring-up status is not a result.
% ----------------------------------------------------------------------------

% \subsection{Embodied Bring-Up (Yahboom + Jetson)}
% \begin{table}[t]
\centering
\caption{Embodied Yahboom bring-up status (verified runs on Jetson + 7-DoF arm).}
\label{tab:yahboom_bringup_status}
\begin{tabular}{p{0.52\linewidth}cc}
\toprule
Check & Outcome & Evidence \\
\midrule
Jetson API reachability (\texttt{/status}, \texttt{/snapshot}) & Pass & Both endpoints return valid payloads \\
Safe bridge smoke test (angles + gripper) & Pass & Command ACK success on all checks \\
Random closed-loop execute (tiny deltas) & Pass & 2/2 \texttt{move\_to}, 2/2 gripper ACK success \\
OpenVLA loop before decode fix & Guarded pass & Motion safety guard prevented saturated-action deltas \\
OpenVLA loop after decode fix & Pass & 3/3 \texttt{move\_to}, 3/3 gripper ACK success \\
Autonomous pick/place task success benchmark & Pending & Not yet measured in standardized trial protocol \\
\bottomrule
\end{tabular}
\end{table}

% \textbf{Scope:} This submission prioritizes LIBERO simulation evidence...
% Physical deployment is summarized only as infrastructure bring-up...

% ----------------------------------------------------------------------------
% CUT 7: OpenVLA OOD status table (was "PENDING")
% Reason: Placeholder with no data.
% ----------------------------------------------------------------------------

% \begin{table}[t]
\centering
\caption{\textbf{OpenVLA OOD baseline campaign on hard tasks 8--9} (seeds s42, s43, s44; $20$ episodes/task; OpenVLA-7B + OOD push). Each cell reports success rate (\%) and successes/episodes across the two hard tasks for that seed. The aggregate row pools all three seeds ($n{=}120$ per controller). All three seeds are complete.}
\label{tab:openvla_ood_status}
\begin{tabular}{lcccc}
\toprule
Seed & Vanilla & Latent Stop & MPPI & Steering \\
\midrule
\texttt{s42} & 52.5 (21/40) & 42.5 (17/40) & 52.5 (21/40) & 55.0 (22/40) \\
\texttt{s43} & 52.5 (21/40) & 42.5 (17/40) & 55.0 (22/40) & 45.0 (18/40) \\
\texttt{s44} & 55.0 (22/40) & 45.0 (18/40) & 57.5 (23/40) & 62.5 (25/40) \\
\midrule
Aggregate (s42--s44) & 53.3 (64/120) & 43.3 (52/120) & 55.0 (66/120) & 54.2 (65/120) \\
\bottomrule
\end{tabular}
\end{table}

